%arara: xelatex
\documentclass{ctexart}
\usepackage{amsmath,amssymb,amsthm,bbm,listings,hyperref,xcolor}
\newtheorem{definition}{Definition}
\newtheorem{theorem}{Theorem}
\newtheorem{problem}{Problem}
\newtheorem{solution}{Solution}
\usepackage{listings}
\usepackage[a4paper,left=2cm,right=3cm,top=2.5cm,bottom=2.5cm]{geometry}
\lstset{
  language=Matlab,
  xleftmargin=2em,
  frame=single,
  basicstyle=\ttfamily\small,
  keepspaces=true,
  columns=fullflexible,
  showstringspaces=false
}
\everymath{\displaystyle}
\newlength\inlineHeight
\newlength\inlineWidth
\long\def\(#1\){%
  \settoheight{\inlineHeight}{$#1$}%
  \settowidth{\inlineWidth}{$#1$}%
  \ifdim \inlineWidth > 0.5\textwidth%
  $$#1$$%
  \else%
  \ifdim \inlineHeight > 1.5em%
  $$#1$$%
  \else%
  $#1$%
  \fi%
  \fi%
}

\title{Iterative Linear Systems}
\author{
  Yinya Wang\\
  25114020018\\
  \href{mailto:yinyawang25@m.fudan.edu.cn}{yinyawang25@m.fudan.edu.cn}
}
\date{}

\begin{document}
\maketitle
\section{Problem 20.18}

Let
\[
  A=
  \begin{bmatrix}
    0.1000 & 0.5000 & -0.1000\\
    0.4000 & 0.2000 & 0.6000\\
    0.2000 & -0.3000 & 0.4000
  \end{bmatrix}.
\]
Since $A$ is not diagonally dominant, the convergence of the Jacobi and Gauss--Seidel methods is not guaranteed.

\paragraph{(a) Norms of $B_J$ and $B_{GS}$}

The computed norms of $B_J$ are
\[
  \|B_J\|_1 = 5.7500,\qquad
  \|B_J\|_2 = 5.2180,\qquad
  \|B_J\|_\infty = 6.0000.
\]

The computed norms of $B_{GS}$ are
\[
  \|B_{GS}\|_1 = 25.0000,\qquad
  \|B_{GS}\|_2 = 16.3566,\qquad
  \|B_{GS}\|_\infty = 15.0000.
\]

\paragraph{(b) Spectral radii}

\[
  \rho(B_J)=3.1756,\qquad
  \rho(B_{GS})=3.7500.
\]

\paragraph{(c) Convergence analysis}

A stationary iteration converges if and only if
\[
  \rho(B)<1.
\]
Since
\[
  \rho(B_J)>1,\qquad \rho(B_{GS})>1,
\]
neither the Jacobi method nor the Gauss--Seidel method converges.

To verify this numerically, we choose a vector $b=1$ and set the initial guess
\[
  x^{(0)}=0.
\]
The tolerance is chosen as
\[
  \text{tol} = 10^{-14},
\]
and the maximum number of iterations is set to
\[
  \text{numiter} = 2{,}000{,}000.
\]

Despite allowing a large number of iterations and a reasonable tolerance, both methods fail to converge:
\[
  \text{GS does not converge},\qquad
  \text{Jacobi does not converge}.
\]

This agrees with the theoretical prediction, since the spectral radii of both iteration matrices are greater than $1$.

\paragraph{Conclusion}

Both the Jacobi and Gauss--Seidel methods do not converge for this matrix. This is consistent with the fact that $A$ is not diagonally dominant and that $\rho(B_J)>1$ and $\rho(B_{GS})>1$.
\paragraph{Codes}
\subsection{Codes for problem 20.18}

{\noindent\hspace*{0ex}\textcolor[HTML]{e05661}{\ttfamily{A}}\ttfamily{ }\textcolor[HTML]{56b6c2}{\ttfamily{=}}\ttfamily{ }\textcolor[HTML]{9a77cf}{\ttfamily{[}}\textcolor[HTML]{ee9025}{\ttfamily{0.1}}\textcolor[HTML]{6a6a6a}{\ttfamily{,}}\ttfamily{ }\textcolor[HTML]{ee9025}{\ttfamily{0.5}}\textcolor[HTML]{6a6a6a}{\ttfamily{,}}\ttfamily{ }\textcolor[HTML]{56b6c2}{\ttfamily{-}}\textcolor[HTML]{ee9025}{\ttfamily{0.1}}\textcolor[HTML]{6a6a6a}{\ttfamily{;}}\ttfamily{ }\textcolor[HTML]{ee9025}{\ttfamily{0.4}}\textcolor[HTML]{6a6a6a}{\ttfamily{,}}\ttfamily{ }\textcolor[HTML]{ee9025}{\ttfamily{0.2}}\textcolor[HTML]{6a6a6a}{\ttfamily{,}}\ttfamily{ }\textcolor[HTML]{ee9025}{\ttfamily{0.6}}\textcolor[HTML]{6a6a6a}{\ttfamily{;}}\ttfamily{ }\textcolor[HTML]{ee9025}{\ttfamily{0.2}}\textcolor[HTML]{6a6a6a}{\ttfamily{,}}\ttfamily{ }\textcolor[HTML]{56b6c2}{\ttfamily{-}}\textcolor[HTML]{ee9025}{\ttfamily{0.3}}\textcolor[HTML]{6a6a6a}{\ttfamily{,}}\ttfamily{ }\textcolor[HTML]{ee9025}{\ttfamily{0.4}}\textcolor[HTML]{9a77cf}{\ttfamily{]}}\textcolor[HTML]{6a6a6a}{\ttfamily{;}}\newline
\hspace*{0ex}\textcolor[HTML]{e05661}{\ttfamily{sizeA}}\ttfamily{ }\textcolor[HTML]{56b6c2}{\ttfamily{=}}\ttfamily{ }\textcolor[HTML]{118dc3}{\ttfamily{size}}\textcolor[HTML]{9a77cf}{\ttfamily{(}}\textcolor[HTML]{e05661}{\ttfamily{A}}\textcolor[HTML]{6a6a6a}{\ttfamily{,}}\ttfamily{ }\textcolor[HTML]{ee9025}{\ttfamily{1}}\textcolor[HTML]{9a77cf}{\ttfamily{)}}\textcolor[HTML]{6a6a6a}{\ttfamily{;}}\newline
\hspace*{0ex}\newline
\hspace*{0ex}\textcolor[HTML]{e05661}{\ttfamily{U}}\ttfamily{ }\textcolor[HTML]{56b6c2}{\ttfamily{=}}\ttfamily{ }\textcolor[HTML]{118dc3}{\ttfamily{triu}}\textcolor[HTML]{9a77cf}{\ttfamily{(}}\textcolor[HTML]{e05661}{\ttfamily{A}}\textcolor[HTML]{6a6a6a}{\ttfamily{,}}\ttfamily{ }\textcolor[HTML]{ee9025}{\ttfamily{1}}\textcolor[HTML]{9a77cf}{\ttfamily{)}}\textcolor[HTML]{6a6a6a}{\ttfamily{;}}\newline
\hspace*{0ex}\textcolor[HTML]{e05661}{\ttfamily{D}}\ttfamily{ }\textcolor[HTML]{56b6c2}{\ttfamily{=}}\ttfamily{ }\textcolor[HTML]{118dc3}{\ttfamily{diag}}\textcolor[HTML]{9a77cf}{\ttfamily{(}}\textcolor[HTML]{118dc3}{\ttfamily{diag}}\textcolor[HTML]{9a77cf}{\ttfamily{(}}\textcolor[HTML]{e05661}{\ttfamily{A}}\textcolor[HTML]{9a77cf}{\ttfamily{))}}\textcolor[HTML]{6a6a6a}{\ttfamily{;}}\newline
\hspace*{0ex}\textcolor[HTML]{e05661}{\ttfamily{L}}\ttfamily{ }\textcolor[HTML]{56b6c2}{\ttfamily{=}}\ttfamily{ }\textcolor[HTML]{118dc3}{\ttfamily{tril}}\textcolor[HTML]{9a77cf}{\ttfamily{(}}\textcolor[HTML]{e05661}{\ttfamily{A}}\textcolor[HTML]{6a6a6a}{\ttfamily{,}}\ttfamily{ }\textcolor[HTML]{56b6c2}{\ttfamily{-}}\textcolor[HTML]{ee9025}{\ttfamily{1}}\textcolor[HTML]{9a77cf}{\ttfamily{)}}\textcolor[HTML]{6a6a6a}{\ttfamily{;}}\newline
\hspace*{0ex}\newline
\hspace*{0ex}\textcolor[HTML]{eea825}{\ttfamily{BGS}}\ttfamily{ }\textcolor[HTML]{56b6c2}{\ttfamily{=}}\ttfamily{ }\textcolor[HTML]{56b6c2}{\ttfamily{-}}\textcolor[HTML]{118dc3}{\ttfamily{inv}}\textcolor[HTML]{9a77cf}{\ttfamily{(}}\textcolor[HTML]{e05661}{\ttfamily{L}}\ttfamily{ }\textcolor[HTML]{56b6c2}{\ttfamily{+}}\ttfamily{ }\textcolor[HTML]{e05661}{\ttfamily{D}}\textcolor[HTML]{9a77cf}{\ttfamily{)}}\ttfamily{ }\textcolor[HTML]{56b6c2}{\ttfamily{*}}\ttfamily{ }\textcolor[HTML]{e05661}{\ttfamily{U}}\textcolor[HTML]{6a6a6a}{\ttfamily{;}}\newline
\hspace*{0ex}\textcolor[HTML]{e05661}{\ttfamily{rhoGS}}\ttfamily{ }\textcolor[HTML]{56b6c2}{\ttfamily{=}}\ttfamily{ }\textcolor[HTML]{118dc3}{\ttfamily{max}}\textcolor[HTML]{9a77cf}{\ttfamily{(}}\textcolor[HTML]{118dc3}{\ttfamily{abs}}\textcolor[HTML]{9a77cf}{\ttfamily{(}}\textcolor[HTML]{118dc3}{\ttfamily{eig}}\textcolor[HTML]{9a77cf}{\ttfamily{(}}\textcolor[HTML]{eea825}{\ttfamily{BGS}}\textcolor[HTML]{9a77cf}{\ttfamily{)))}}\textcolor[HTML]{6a6a6a}{\ttfamily{;}}\newline
\hspace*{0ex}\textcolor[HTML]{e05661}{\ttfamily{n1\_GS}}\ttfamily{ }\textcolor[HTML]{56b6c2}{\ttfamily{=}}\ttfamily{ }\textcolor[HTML]{118dc3}{\ttfamily{norm}}\textcolor[HTML]{9a77cf}{\ttfamily{(}}\textcolor[HTML]{eea825}{\ttfamily{BGS}}\textcolor[HTML]{6a6a6a}{\ttfamily{,}}\ttfamily{ }\textcolor[HTML]{ee9025}{\ttfamily{1}}\textcolor[HTML]{9a77cf}{\ttfamily{)}}\textcolor[HTML]{6a6a6a}{\ttfamily{;}}\newline
\hspace*{0ex}\textcolor[HTML]{e05661}{\ttfamily{n2\_GS}}\ttfamily{ }\textcolor[HTML]{56b6c2}{\ttfamily{=}}\ttfamily{ }\textcolor[HTML]{118dc3}{\ttfamily{norm}}\textcolor[HTML]{9a77cf}{\ttfamily{(}}\textcolor[HTML]{eea825}{\ttfamily{BGS}}\textcolor[HTML]{6a6a6a}{\ttfamily{,}}\ttfamily{ }\textcolor[HTML]{ee9025}{\ttfamily{2}}\textcolor[HTML]{9a77cf}{\ttfamily{)}}\textcolor[HTML]{6a6a6a}{\ttfamily{;}}\newline
\hspace*{0ex}\textcolor[HTML]{e05661}{\ttfamily{nin\_GS}}\ttfamily{ }\textcolor[HTML]{56b6c2}{\ttfamily{=}}\ttfamily{ }\textcolor[HTML]{118dc3}{\ttfamily{norm}}\textcolor[HTML]{9a77cf}{\ttfamily{(}}\textcolor[HTML]{eea825}{\ttfamily{BGS}}\textcolor[HTML]{6a6a6a}{\ttfamily{,}}\ttfamily{ }\textcolor[HTML]{e05661}{\ttfamily{inf}}\textcolor[HTML]{9a77cf}{\ttfamily{)}}\textcolor[HTML]{6a6a6a}{\ttfamily{;}}\newline
\hspace*{0ex}\textcolor[HTML]{118dc3}{\ttfamily{fprintf}}\textcolor[HTML]{9a77cf}{\ttfamily{(}}\textcolor[HTML]{56b6c2}{\ttfamily{'}}\textcolor[HTML]{1da912}{\ttfamily{1-norm of BGS = }}\textcolor[HTML]{118dc3}{\ttfamily{\%.4f}}\textcolor[HTML]{56b6c2}{\ttfamily{\textbackslash{}n}}\textcolor[HTML]{56b6c2}{\ttfamily{'}}\textcolor[HTML]{6a6a6a}{\ttfamily{,}}\ttfamily{ }\textcolor[HTML]{e05661}{\ttfamily{n1\_GS}}\textcolor[HTML]{9a77cf}{\ttfamily{)}}\textcolor[HTML]{6a6a6a}{\ttfamily{;}}\newline
\hspace*{0ex}\textcolor[HTML]{118dc3}{\ttfamily{fprintf}}\textcolor[HTML]{9a77cf}{\ttfamily{(}}\textcolor[HTML]{56b6c2}{\ttfamily{'}}\textcolor[HTML]{1da912}{\ttfamily{2-norm of BGS = }}\textcolor[HTML]{118dc3}{\ttfamily{\%.4f}}\textcolor[HTML]{56b6c2}{\ttfamily{\textbackslash{}n}}\textcolor[HTML]{56b6c2}{\ttfamily{'}}\textcolor[HTML]{6a6a6a}{\ttfamily{,}}\ttfamily{ }\textcolor[HTML]{e05661}{\ttfamily{n2\_GS}}\textcolor[HTML]{9a77cf}{\ttfamily{)}}\textcolor[HTML]{6a6a6a}{\ttfamily{;}}\newline
\hspace*{0ex}\textcolor[HTML]{118dc3}{\ttfamily{fprintf}}\textcolor[HTML]{9a77cf}{\ttfamily{(}}\textcolor[HTML]{56b6c2}{\ttfamily{'}}\textcolor[HTML]{1da912}{\ttfamily{infinity-norm of BGS = }}\textcolor[HTML]{118dc3}{\ttfamily{\%.4f}}\textcolor[HTML]{56b6c2}{\ttfamily{\textbackslash{}n}}\textcolor[HTML]{56b6c2}{\ttfamily{'}}\textcolor[HTML]{6a6a6a}{\ttfamily{,}}\ttfamily{ }\textcolor[HTML]{e05661}{\ttfamily{nin\_GS}}\textcolor[HTML]{9a77cf}{\ttfamily{)}}\textcolor[HTML]{6a6a6a}{\ttfamily{;}}\newline
\hspace*{0ex}\textcolor[HTML]{118dc3}{\ttfamily{fprintf}}\textcolor[HTML]{9a77cf}{\ttfamily{(}}\textcolor[HTML]{56b6c2}{\ttfamily{'}}\textcolor[HTML]{1da912}{\ttfamily{Spectral radius of BGS = }}\textcolor[HTML]{118dc3}{\ttfamily{\%.4f}}\textcolor[HTML]{56b6c2}{\ttfamily{\textbackslash{}n}}\textcolor[HTML]{56b6c2}{\ttfamily{'}}\textcolor[HTML]{6a6a6a}{\ttfamily{,}}\ttfamily{ }\textcolor[HTML]{e05661}{\ttfamily{rhoGS}}\textcolor[HTML]{9a77cf}{\ttfamily{)}}\textcolor[HTML]{6a6a6a}{\ttfamily{;}}\newline
\hspace*{0ex}\newline
\hspace*{0ex}\textcolor[HTML]{eea825}{\ttfamily{BJ}}\ttfamily{ }\textcolor[HTML]{56b6c2}{\ttfamily{=}}\ttfamily{ }\textcolor[HTML]{56b6c2}{\ttfamily{-}}\textcolor[HTML]{118dc3}{\ttfamily{inv}}\textcolor[HTML]{9a77cf}{\ttfamily{(}}\textcolor[HTML]{e05661}{\ttfamily{D}}\textcolor[HTML]{9a77cf}{\ttfamily{)}}\ttfamily{ }\textcolor[HTML]{56b6c2}{\ttfamily{*}}\ttfamily{ }\textcolor[HTML]{9a77cf}{\ttfamily{(}}\textcolor[HTML]{e05661}{\ttfamily{L}}\ttfamily{ }\textcolor[HTML]{56b6c2}{\ttfamily{+}}\ttfamily{ }\textcolor[HTML]{e05661}{\ttfamily{U}}\textcolor[HTML]{9a77cf}{\ttfamily{)}}\textcolor[HTML]{6a6a6a}{\ttfamily{;}}\newline
\hspace*{0ex}\textcolor[HTML]{e05661}{\ttfamily{n1\_J}}\ttfamily{ }\textcolor[HTML]{56b6c2}{\ttfamily{=}}\ttfamily{ }\textcolor[HTML]{118dc3}{\ttfamily{norm}}\textcolor[HTML]{9a77cf}{\ttfamily{(}}\textcolor[HTML]{eea825}{\ttfamily{BJ}}\textcolor[HTML]{6a6a6a}{\ttfamily{,}}\ttfamily{ }\textcolor[HTML]{ee9025}{\ttfamily{1}}\textcolor[HTML]{9a77cf}{\ttfamily{)}}\textcolor[HTML]{6a6a6a}{\ttfamily{;}}\newline
\hspace*{0ex}\textcolor[HTML]{e05661}{\ttfamily{n2\_J}}\ttfamily{ }\textcolor[HTML]{56b6c2}{\ttfamily{=}}\ttfamily{ }\textcolor[HTML]{118dc3}{\ttfamily{norm}}\textcolor[HTML]{9a77cf}{\ttfamily{(}}\textcolor[HTML]{eea825}{\ttfamily{BJ}}\textcolor[HTML]{6a6a6a}{\ttfamily{,}}\ttfamily{ }\textcolor[HTML]{ee9025}{\ttfamily{2}}\textcolor[HTML]{9a77cf}{\ttfamily{)}}\textcolor[HTML]{6a6a6a}{\ttfamily{;}}\newline
\hspace*{0ex}\textcolor[HTML]{e05661}{\ttfamily{nin\_J}}\ttfamily{ }\textcolor[HTML]{56b6c2}{\ttfamily{=}}\ttfamily{ }\textcolor[HTML]{118dc3}{\ttfamily{norm}}\textcolor[HTML]{9a77cf}{\ttfamily{(}}\textcolor[HTML]{eea825}{\ttfamily{BJ}}\textcolor[HTML]{6a6a6a}{\ttfamily{,}}\ttfamily{ }\textcolor[HTML]{e05661}{\ttfamily{inf}}\textcolor[HTML]{9a77cf}{\ttfamily{)}}\textcolor[HTML]{6a6a6a}{\ttfamily{;}}\newline
\hspace*{0ex}\textcolor[HTML]{e05661}{\ttfamily{rhoBJ}}\ttfamily{ }\textcolor[HTML]{56b6c2}{\ttfamily{=}}\ttfamily{ }\textcolor[HTML]{118dc3}{\ttfamily{max}}\textcolor[HTML]{9a77cf}{\ttfamily{(}}\textcolor[HTML]{118dc3}{\ttfamily{abs}}\textcolor[HTML]{9a77cf}{\ttfamily{(}}\textcolor[HTML]{118dc3}{\ttfamily{eig}}\textcolor[HTML]{9a77cf}{\ttfamily{(}}\textcolor[HTML]{eea825}{\ttfamily{BJ}}\textcolor[HTML]{9a77cf}{\ttfamily{)))}}\textcolor[HTML]{6a6a6a}{\ttfamily{;}}\newline
\hspace*{0ex}\textcolor[HTML]{118dc3}{\ttfamily{fprintf}}\textcolor[HTML]{9a77cf}{\ttfamily{(}}\textcolor[HTML]{56b6c2}{\ttfamily{'}}\textcolor[HTML]{1da912}{\ttfamily{1-norm of BJ = }}\textcolor[HTML]{118dc3}{\ttfamily{\%.4f}}\textcolor[HTML]{56b6c2}{\ttfamily{\textbackslash{}n}}\textcolor[HTML]{56b6c2}{\ttfamily{'}}\textcolor[HTML]{6a6a6a}{\ttfamily{,}}\ttfamily{ }\textcolor[HTML]{e05661}{\ttfamily{n1\_J}}\textcolor[HTML]{9a77cf}{\ttfamily{)}}\textcolor[HTML]{6a6a6a}{\ttfamily{;}}\newline
\hspace*{0ex}\textcolor[HTML]{118dc3}{\ttfamily{fprintf}}\textcolor[HTML]{9a77cf}{\ttfamily{(}}\textcolor[HTML]{56b6c2}{\ttfamily{'}}\textcolor[HTML]{1da912}{\ttfamily{2-norm of BJ = }}\textcolor[HTML]{118dc3}{\ttfamily{\%.4f}}\textcolor[HTML]{56b6c2}{\ttfamily{\textbackslash{}n}}\textcolor[HTML]{56b6c2}{\ttfamily{'}}\textcolor[HTML]{6a6a6a}{\ttfamily{,}}\ttfamily{ }\textcolor[HTML]{e05661}{\ttfamily{n2\_J}}\textcolor[HTML]{9a77cf}{\ttfamily{)}}\textcolor[HTML]{6a6a6a}{\ttfamily{;}}\newline
\hspace*{0ex}\textcolor[HTML]{118dc3}{\ttfamily{fprintf}}\textcolor[HTML]{9a77cf}{\ttfamily{(}}\textcolor[HTML]{56b6c2}{\ttfamily{'}}\textcolor[HTML]{1da912}{\ttfamily{infinity-norm of BJ = }}\textcolor[HTML]{118dc3}{\ttfamily{\%.4f}}\textcolor[HTML]{56b6c2}{\ttfamily{\textbackslash{}n}}\textcolor[HTML]{56b6c2}{\ttfamily{'}}\textcolor[HTML]{6a6a6a}{\ttfamily{,}}\ttfamily{ }\textcolor[HTML]{e05661}{\ttfamily{nin\_J}}\textcolor[HTML]{9a77cf}{\ttfamily{)}}\textcolor[HTML]{6a6a6a}{\ttfamily{;}}\newline
\hspace*{0ex}\textcolor[HTML]{118dc3}{\ttfamily{fprintf}}\textcolor[HTML]{9a77cf}{\ttfamily{(}}\textcolor[HTML]{56b6c2}{\ttfamily{'}}\textcolor[HTML]{1da912}{\ttfamily{Spectral radius of BJ = }}\textcolor[HTML]{118dc3}{\ttfamily{\%.4f}}\textcolor[HTML]{56b6c2}{\ttfamily{\textbackslash{}n}}\textcolor[HTML]{56b6c2}{\ttfamily{'}}\textcolor[HTML]{6a6a6a}{\ttfamily{,}}\ttfamily{ }\textcolor[HTML]{e05661}{\ttfamily{rhoBJ}}\textcolor[HTML]{9a77cf}{\ttfamily{)}}\textcolor[HTML]{6a6a6a}{\ttfamily{;}}\newline
\hspace*{0ex}\newline
\hspace*{0ex}\textcolor[HTML]{e05661}{\ttfamily{tol}}\ttfamily{ }\textcolor[HTML]{56b6c2}{\ttfamily{=}}\ttfamily{ }\textcolor[HTML]{ee9025}{\ttfamily{1.0e-14}}\textcolor[HTML]{6a6a6a}{\ttfamily{;}}\newline
\hspace*{0ex}\textcolor[HTML]{e05661}{\ttfamily{numiter}}\ttfamily{ }\textcolor[HTML]{56b6c2}{\ttfamily{=}}\ttfamily{ }\textcolor[HTML]{ee9025}{\ttfamily{2000000}}\textcolor[HTML]{6a6a6a}{\ttfamily{;}}\newline
\hspace*{0ex}\textcolor[HTML]{e05661}{\ttfamily{b}}\ttfamily{ }\textcolor[HTML]{56b6c2}{\ttfamily{=}}\ttfamily{ }\textcolor[HTML]{118dc3}{\ttfamily{ones}}\textcolor[HTML]{9a77cf}{\ttfamily{(}}\textcolor[HTML]{e05661}{\ttfamily{sizeA}}\textcolor[HTML]{6a6a6a}{\ttfamily{,}}\ttfamily{ }\textcolor[HTML]{ee9025}{\ttfamily{1}}\textcolor[HTML]{9a77cf}{\ttfamily{)}}\textcolor[HTML]{6a6a6a}{\ttfamily{;}}\newline
\hspace*{0ex}\textcolor[HTML]{e05661}{\ttfamily{x0}}\ttfamily{ }\textcolor[HTML]{56b6c2}{\ttfamily{=}}\ttfamily{ }\textcolor[HTML]{118dc3}{\ttfamily{zeros}}\textcolor[HTML]{9a77cf}{\ttfamily{(}}\textcolor[HTML]{e05661}{\ttfamily{sizeA}}\textcolor[HTML]{6a6a6a}{\ttfamily{,}}\ttfamily{ }\textcolor[HTML]{ee9025}{\ttfamily{1}}\textcolor[HTML]{9a77cf}{\ttfamily{)}}\textcolor[HTML]{6a6a6a}{\ttfamily{;}}\newline
\hspace*{0ex}\textcolor[HTML]{9a77cf}{\ttfamily{[}}\textcolor[HTML]{9a77cf}{\ttfamily{\textasciitilde{}}}\textcolor[HTML]{6a6a6a}{\ttfamily{,}}\ttfamily{ }\textcolor[HTML]{eea825}{\ttfamily{GSiter}}\textcolor[HTML]{6a6a6a}{\ttfamily{,}}\ttfamily{ }\textcolor[HTML]{eea825}{\ttfamily{GSresider}}\textcolor[HTML]{9a77cf}{\ttfamily{]}}\ttfamily{ }\textcolor[HTML]{56b6c2}{\ttfamily{=}}\ttfamily{ }\textcolor[HTML]{118dc3}{\ttfamily{gausseidel}}\textcolor[HTML]{9a77cf}{\ttfamily{(}}\textcolor[HTML]{e05661}{\ttfamily{A}}\textcolor[HTML]{6a6a6a}{\ttfamily{,}}\ttfamily{ }\textcolor[HTML]{e05661}{\ttfamily{b}}\textcolor[HTML]{6a6a6a}{\ttfamily{,}}\ttfamily{ }\textcolor[HTML]{e05661}{\ttfamily{x0}}\textcolor[HTML]{6a6a6a}{\ttfamily{,}}\ttfamily{ }\textcolor[HTML]{e05661}{\ttfamily{tol}}\textcolor[HTML]{6a6a6a}{\ttfamily{,}}\ttfamily{ }\textcolor[HTML]{e05661}{\ttfamily{numiter}}\textcolor[HTML]{9a77cf}{\ttfamily{)}}\textcolor[HTML]{6a6a6a}{\ttfamily{;}}\newline
\hspace*{0ex}\textcolor[HTML]{9a77cf}{\ttfamily{if}}\ttfamily{ }\textcolor[HTML]{eea825}{\ttfamily{GSiter}}\ttfamily{ }\textcolor[HTML]{56b6c2}{\ttfamily{>}}\ttfamily{ }\textcolor[HTML]{ee9025}{\ttfamily{0}}\newline
\hspace*{2ex}\textcolor[HTML]{118dc3}{\ttfamily{fprintf}}\textcolor[HTML]{9a77cf}{\ttfamily{(}}\textcolor[HTML]{56b6c2}{\ttfamily{'}}\textcolor[HTML]{1da912}{\ttfamily{Gauss-Seidel converges at step }}\textcolor[HTML]{118dc3}{\ttfamily{\%.4f}}\textcolor[HTML]{1da912}{\ttfamily{ with residual }}\textcolor[HTML]{118dc3}{\ttfamily{\%.4f}}\textcolor[HTML]{56b6c2}{\ttfamily{\textbackslash{}n}}\textcolor[HTML]{56b6c2}{\ttfamily{'}}\textcolor[HTML]{6a6a6a}{\ttfamily{,}}\ttfamily{ }\textcolor[HTML]{9b9fa6}{\ttfamily{...}}\newline
\hspace*{10ex}\textcolor[HTML]{eea825}{\ttfamily{GSiter}}\textcolor[HTML]{6a6a6a}{\ttfamily{,}}\ttfamily{ }\textcolor[HTML]{eea825}{\ttfamily{GSresider}}\textcolor[HTML]{9a77cf}{\ttfamily{)}}\textcolor[HTML]{6a6a6a}{\ttfamily{;}}\newline
\hspace*{0ex}\textcolor[HTML]{9a77cf}{\ttfamily{else}}\newline
\hspace*{2ex}\textcolor[HTML]{118dc3}{\ttfamily{fprintf}}\textcolor[HTML]{9a77cf}{\ttfamily{(}}\textcolor[HTML]{56b6c2}{\ttfamily{'}}\textcolor[HTML]{1da912}{\ttfamily{Gauss-Seidel does not converge}}\textcolor[HTML]{56b6c2}{\ttfamily{\textbackslash{}n}}\textcolor[HTML]{56b6c2}{\ttfamily{'}}\textcolor[HTML]{9a77cf}{\ttfamily{)}}\textcolor[HTML]{6a6a6a}{\ttfamily{;}}\newline
\hspace*{0ex}\textcolor[HTML]{9a77cf}{\ttfamily{end}}\newline
\hspace*{0ex}\textcolor[HTML]{9a77cf}{\ttfamily{[}}\textcolor[HTML]{9a77cf}{\ttfamily{\textasciitilde{}}}\textcolor[HTML]{6a6a6a}{\ttfamily{,}}\ttfamily{ }\textcolor[HTML]{eea825}{\ttfamily{Jiter}}\textcolor[HTML]{6a6a6a}{\ttfamily{,}}\ttfamily{ }\textcolor[HTML]{eea825}{\ttfamily{Jresider}}\textcolor[HTML]{9a77cf}{\ttfamily{]}}\ttfamily{ }\textcolor[HTML]{56b6c2}{\ttfamily{=}}\ttfamily{ }\textcolor[HTML]{118dc3}{\ttfamily{Jacobi}}\textcolor[HTML]{9a77cf}{\ttfamily{(}}\textcolor[HTML]{e05661}{\ttfamily{A}}\textcolor[HTML]{6a6a6a}{\ttfamily{,}}\ttfamily{ }\textcolor[HTML]{e05661}{\ttfamily{b}}\textcolor[HTML]{6a6a6a}{\ttfamily{,}}\ttfamily{ }\textcolor[HTML]{e05661}{\ttfamily{x0}}\textcolor[HTML]{6a6a6a}{\ttfamily{,}}\ttfamily{ }\textcolor[HTML]{e05661}{\ttfamily{tol}}\textcolor[HTML]{6a6a6a}{\ttfamily{,}}\ttfamily{ }\textcolor[HTML]{e05661}{\ttfamily{numiter}}\textcolor[HTML]{9a77cf}{\ttfamily{)}}\textcolor[HTML]{6a6a6a}{\ttfamily{;}}\newline
\hspace*{0ex}\textcolor[HTML]{9a77cf}{\ttfamily{if}}\ttfamily{ }\textcolor[HTML]{eea825}{\ttfamily{Jiter}}\ttfamily{ }\textcolor[HTML]{56b6c2}{\ttfamily{>}}\ttfamily{ }\textcolor[HTML]{ee9025}{\ttfamily{0}}\newline
\hspace*{2ex}\textcolor[HTML]{118dc3}{\ttfamily{fprintf}}\textcolor[HTML]{9a77cf}{\ttfamily{(}}\textcolor[HTML]{56b6c2}{\ttfamily{'}}\textcolor[HTML]{1da912}{\ttfamily{Jacobi converges at step }}\textcolor[HTML]{118dc3}{\ttfamily{\%.4f}}\textcolor[HTML]{1da912}{\ttfamily{ with residual }}\textcolor[HTML]{118dc3}{\ttfamily{\%.4f}}\textcolor[HTML]{56b6c2}{\ttfamily{\textbackslash{}n}}\textcolor[HTML]{56b6c2}{\ttfamily{'}}\textcolor[HTML]{6a6a6a}{\ttfamily{,}}\ttfamily{ }\textcolor[HTML]{9b9fa6}{\ttfamily{...}}\newline
\hspace*{10ex}\textcolor[HTML]{eea825}{\ttfamily{Jiter}}\textcolor[HTML]{6a6a6a}{\ttfamily{,}}\ttfamily{ }\textcolor[HTML]{eea825}{\ttfamily{Jresider}}\textcolor[HTML]{9a77cf}{\ttfamily{)}}\textcolor[HTML]{6a6a6a}{\ttfamily{;}}\newline
\hspace*{0ex}\textcolor[HTML]{9a77cf}{\ttfamily{else}}\newline
\hspace*{2ex}\textcolor[HTML]{118dc3}{\ttfamily{fprintf}}\textcolor[HTML]{9a77cf}{\ttfamily{(}}\textcolor[HTML]{56b6c2}{\ttfamily{'}}\textcolor[HTML]{1da912}{\ttfamily{Jacobi does not converge}}\textcolor[HTML]{56b6c2}{\ttfamily{\textbackslash{}n}}\textcolor[HTML]{56b6c2}{\ttfamily{'}}\textcolor[HTML]{9a77cf}{\ttfamily{)}}\textcolor[HTML]{6a6a6a}{\ttfamily{;}}\newline
\hspace*{0ex}\textcolor[HTML]{9a77cf}{\ttfamily{end}}
}
% \subsection{Codes for Jacobi}
% {\noindent\input{Jacobi.tex}}
% \subsection{Codes for GS}
% {\noindent\hspace*{0ex}\textcolor[HTML]{9a77cf}{\ttfamily{function}}\ttfamily{ }\textcolor[HTML]{9a77cf}{\ttfamily{[}}\textcolor[HTML]{e05661}{\ttfamily{x}}\textcolor[HTML]{6a6a6a}{\ttfamily{,}}\ttfamily{ }\textcolor[HTML]{e05661}{\ttfamily{iter}}\textcolor[HTML]{6a6a6a}{\ttfamily{,}}\ttfamily{ }\textcolor[HTML]{e05661}{\ttfamily{relresid}}\textcolor[HTML]{9a77cf}{\ttfamily{]}}\ttfamily{ }\textcolor[HTML]{56b6c2}{\ttfamily{=}}\ttfamily{ }\textcolor[HTML]{118dc3}{\ttfamily{GS}}\textcolor[HTML]{9a77cf}{\ttfamily{(}}\textcolor[HTML]{e05661}{\ttfamily{A}}\textcolor[HTML]{6a6a6a}{\ttfamily{,}}\ttfamily{ }\textcolor[HTML]{e05661}{\ttfamily{b}}\textcolor[HTML]{6a6a6a}{\ttfamily{,}}\ttfamily{ }\textcolor[HTML]{e05661}{\ttfamily{x0}}\textcolor[HTML]{6a6a6a}{\ttfamily{,}}\ttfamily{ }\textcolor[HTML]{e05661}{\ttfamily{tol}}\textcolor[HTML]{6a6a6a}{\ttfamily{,}}\ttfamily{ }\textcolor[HTML]{e05661}{\ttfamily{maxiter}}\textcolor[HTML]{9a77cf}{\ttfamily{)}}\newline
\hspace*{2ex}\textcolor[HTML]{9b9fa6}{\ttfamily{\% GS Successive overrelaxation iteration.}}\newline
\hspace*{2ex}\textcolor[HTML]{9b9fa6}{\ttfamily{\%}}\newline
\hspace*{2ex}\textcolor[HTML]{9b9fa6}{\ttfamily{\%   [x,iter,relresid] = sor(A,b,x0,tol,maxiter) computes}}\newline
\hspace*{2ex}\textcolor[HTML]{9b9fa6}{\ttfamily{\%   the solution x of Ax = b using the GS iterative method.}}\newline
\hspace*{2ex}\textcolor[HTML]{9b9fa6}{\ttfamily{\%   x0 is the initial approximation to the solution,}}\newline
\hspace*{2ex}\textcolor[HTML]{9b9fa6}{\ttfamily{\%   is the relaxation parameter, tol is the error tolerance,}}\newline
\hspace*{2ex}\textcolor[HTML]{9b9fa6}{\ttfamily{\%   and maxiter is the number of iterations to execute.}}\newline
\hspace*{2ex}\textcolor[HTML]{9b9fa6}{\ttfamily{\%   If the GS method converges, iter contains the number of}}\newline
\hspace*{2ex}\textcolor[HTML]{9b9fa6}{\ttfamily{\%   iterations required to meet the error tolerance.}}\newline
\hspace*{2ex}\textcolor[HTML]{9b9fa6}{\ttfamily{\%   If the GS iteration does not converge, iter = -1.}}\newline
\hspace*{2ex}\textcolor[HTML]{9b9fa6}{\ttfamily{\%   relresid is the relative residual obtained by the iteration.}}\newline
\hspace*{0ex}\newline
\hspace*{2ex}\textcolor[HTML]{9a77cf}{\ttfamily{[}}\textcolor[HTML]{e05661}{\ttfamily{m}}\textcolor[HTML]{6a6a6a}{\ttfamily{,}}\ttfamily{ }\textcolor[HTML]{e05661}{\ttfamily{n}}\textcolor[HTML]{9a77cf}{\ttfamily{]}}\ttfamily{ }\textcolor[HTML]{56b6c2}{\ttfamily{=}}\ttfamily{ }\textcolor[HTML]{118dc3}{\ttfamily{size}}\textcolor[HTML]{9a77cf}{\ttfamily{(}}\textcolor[HTML]{e05661}{\ttfamily{A}}\textcolor[HTML]{9a77cf}{\ttfamily{)}}\textcolor[HTML]{6a6a6a}{\ttfamily{;}}\newline
\hspace*{2ex}\textcolor[HTML]{9a77cf}{\ttfamily{if}}\ttfamily{ }\textcolor[HTML]{e05661}{\ttfamily{m}}\ttfamily{ }\textcolor[HTML]{56b6c2}{\ttfamily{\textasciitilde{}=}}\ttfamily{ }\textcolor[HTML]{e05661}{\ttfamily{n}}\newline
\hspace*{4ex}\textcolor[HTML]{118dc3}{\ttfamily{error}}\textcolor[HTML]{9a77cf}{\ttfamily{(}}\textcolor[HTML]{56b6c2}{\ttfamily{'}}\textcolor[HTML]{1da912}{\ttfamily{matrix A  is not square}}\textcolor[HTML]{56b6c2}{\ttfamily{'}}\textcolor[HTML]{9a77cf}{\ttfamily{)}}\textcolor[HTML]{6a6a6a}{\ttfamily{;}}\newline
\hspace*{2ex}\textcolor[HTML]{9a77cf}{\ttfamily{end}}\newline
\hspace*{0ex}\newline
\hspace*{2ex}\textcolor[HTML]{e05661}{\ttfamily{k}}\ttfamily{ }\textcolor[HTML]{56b6c2}{\ttfamily{=}}\ttfamily{ }\textcolor[HTML]{ee9025}{\ttfamily{1}}\textcolor[HTML]{6a6a6a}{\ttfamily{;}}\newline
\hspace*{2ex}\textcolor[HTML]{e05661}{\ttfamily{x}}\ttfamily{ }\textcolor[HTML]{56b6c2}{\ttfamily{=}}\ttfamily{ }\textcolor[HTML]{e05661}{\ttfamily{x0}}\textcolor[HTML]{6a6a6a}{\ttfamily{;}}\newline
\hspace*{2ex}\textcolor[HTML]{9a77cf}{\ttfamily{for}}\ttfamily{ }\textcolor[HTML]{e05661}{\ttfamily{k}}\ttfamily{ }\textcolor[HTML]{56b6c2}{\ttfamily{=}}\ttfamily{ }\textcolor[HTML]{ee9025}{\ttfamily{1}}\textcolor[HTML]{56b6c2}{\ttfamily{:}}\textcolor[HTML]{e05661}{\ttfamily{maxiter}}\newline
\hspace*{4ex}\textcolor[HTML]{118dc3}{\ttfamily{x}}\textcolor[HTML]{9a77cf}{\ttfamily{(}}\textcolor[HTML]{ee9025}{\ttfamily{1}}\textcolor[HTML]{9a77cf}{\ttfamily{)}}\ttfamily{ }\textcolor[HTML]{56b6c2}{\ttfamily{=}}\ttfamily{ }\textcolor[HTML]{ee9025}{\ttfamily{1.0}}\ttfamily{ }\textcolor[HTML]{56b6c2}{\ttfamily{/}}\ttfamily{ }\textcolor[HTML]{118dc3}{\ttfamily{A}}\textcolor[HTML]{9a77cf}{\ttfamily{(}}\textcolor[HTML]{ee9025}{\ttfamily{1}}\textcolor[HTML]{6a6a6a}{\ttfamily{,}}\ttfamily{ }\textcolor[HTML]{ee9025}{\ttfamily{1}}\textcolor[HTML]{9a77cf}{\ttfamily{)}}\ttfamily{ }\textcolor[HTML]{56b6c2}{\ttfamily{*}}\ttfamily{ }\textcolor[HTML]{9a77cf}{\ttfamily{(}}\textcolor[HTML]{118dc3}{\ttfamily{b}}\textcolor[HTML]{9a77cf}{\ttfamily{(}}\textcolor[HTML]{ee9025}{\ttfamily{1}}\textcolor[HTML]{9a77cf}{\ttfamily{)}}\ttfamily{ }\textcolor[HTML]{56b6c2}{\ttfamily{-}}\ttfamily{ }\textcolor[HTML]{118dc3}{\ttfamily{A}}\textcolor[HTML]{9a77cf}{\ttfamily{(}}\textcolor[HTML]{ee9025}{\ttfamily{1}}\textcolor[HTML]{6a6a6a}{\ttfamily{,}}\ttfamily{ }\textcolor[HTML]{ee9025}{\ttfamily{2}}\textcolor[HTML]{56b6c2}{\ttfamily{:}}\textcolor[HTML]{e05661}{\ttfamily{n}}\textcolor[HTML]{9a77cf}{\ttfamily{)}}\ttfamily{ }\textcolor[HTML]{56b6c2}{\ttfamily{*}}\ttfamily{ }\textcolor[HTML]{118dc3}{\ttfamily{x}}\textcolor[HTML]{9a77cf}{\ttfamily{(}}\textcolor[HTML]{ee9025}{\ttfamily{2}}\textcolor[HTML]{56b6c2}{\ttfamily{:}}\textcolor[HTML]{e05661}{\ttfamily{n}}\textcolor[HTML]{9a77cf}{\ttfamily{))}}\textcolor[HTML]{6a6a6a}{\ttfamily{;}}\newline
\hspace*{4ex}\textcolor[HTML]{9a77cf}{\ttfamily{for}}\ttfamily{ }\textcolor[HTML]{e05661}{\ttfamily{i}}\ttfamily{ }\textcolor[HTML]{56b6c2}{\ttfamily{=}}\ttfamily{ }\textcolor[HTML]{ee9025}{\ttfamily{2}}\textcolor[HTML]{56b6c2}{\ttfamily{:}}\textcolor[HTML]{e05661}{\ttfamily{n}}\ttfamily{ }\textcolor[HTML]{56b6c2}{\ttfamily{-}}\ttfamily{ }\textcolor[HTML]{ee9025}{\ttfamily{1}}\newline
\hspace*{6ex}\textcolor[HTML]{118dc3}{\ttfamily{x}}\textcolor[HTML]{9a77cf}{\ttfamily{(}}\textcolor[HTML]{e05661}{\ttfamily{i}}\textcolor[HTML]{9a77cf}{\ttfamily{)}}\ttfamily{ }\textcolor[HTML]{56b6c2}{\ttfamily{=}}\ttfamily{ }\textcolor[HTML]{ee9025}{\ttfamily{1.0}}\ttfamily{ }\textcolor[HTML]{56b6c2}{\ttfamily{/}}\ttfamily{ }\textcolor[HTML]{118dc3}{\ttfamily{A}}\textcolor[HTML]{9a77cf}{\ttfamily{(}}\textcolor[HTML]{e05661}{\ttfamily{i}}\textcolor[HTML]{6a6a6a}{\ttfamily{,}}\ttfamily{ }\textcolor[HTML]{e05661}{\ttfamily{i}}\textcolor[HTML]{9a77cf}{\ttfamily{)}}\ttfamily{ }\textcolor[HTML]{56b6c2}{\ttfamily{*}}\ttfamily{ }\textcolor[HTML]{9a77cf}{\ttfamily{(}}\textcolor[HTML]{118dc3}{\ttfamily{b}}\textcolor[HTML]{9a77cf}{\ttfamily{(}}\textcolor[HTML]{e05661}{\ttfamily{i}}\textcolor[HTML]{9a77cf}{\ttfamily{)}}\ttfamily{ }\textcolor[HTML]{56b6c2}{\ttfamily{-}}\ttfamily{ }\textcolor[HTML]{9b9fa6}{\ttfamily{...}}\newline
\hspace*{30ex}\textcolor[HTML]{118dc3}{\ttfamily{A}}\textcolor[HTML]{9a77cf}{\ttfamily{(}}\textcolor[HTML]{e05661}{\ttfamily{i}}\textcolor[HTML]{6a6a6a}{\ttfamily{,}}\ttfamily{ }\textcolor[HTML]{ee9025}{\ttfamily{1}}\textcolor[HTML]{56b6c2}{\ttfamily{:}}\textcolor[HTML]{e05661}{\ttfamily{i}}\ttfamily{ }\textcolor[HTML]{56b6c2}{\ttfamily{-}}\ttfamily{ }\textcolor[HTML]{ee9025}{\ttfamily{1}}\textcolor[HTML]{9a77cf}{\ttfamily{)}}\ttfamily{ }\textcolor[HTML]{56b6c2}{\ttfamily{*}}\ttfamily{ }\textcolor[HTML]{118dc3}{\ttfamily{x}}\textcolor[HTML]{9a77cf}{\ttfamily{(}}\textcolor[HTML]{ee9025}{\ttfamily{1}}\textcolor[HTML]{56b6c2}{\ttfamily{:}}\textcolor[HTML]{e05661}{\ttfamily{i}}\ttfamily{ }\textcolor[HTML]{56b6c2}{\ttfamily{-}}\ttfamily{ }\textcolor[HTML]{ee9025}{\ttfamily{1}}\textcolor[HTML]{9a77cf}{\ttfamily{)}}\ttfamily{ }\textcolor[HTML]{56b6c2}{\ttfamily{-}}\ttfamily{ }\textcolor[HTML]{118dc3}{\ttfamily{A}}\textcolor[HTML]{9a77cf}{\ttfamily{(}}\textcolor[HTML]{e05661}{\ttfamily{i}}\textcolor[HTML]{6a6a6a}{\ttfamily{,}}\ttfamily{ }\textcolor[HTML]{e05661}{\ttfamily{i}}\ttfamily{ }\textcolor[HTML]{56b6c2}{\ttfamily{+}}\ttfamily{ }\textcolor[HTML]{ee9025}{\ttfamily{1}}\textcolor[HTML]{56b6c2}{\ttfamily{:}}\textcolor[HTML]{e05661}{\ttfamily{n}}\textcolor[HTML]{9a77cf}{\ttfamily{)}}\ttfamily{ }\textcolor[HTML]{56b6c2}{\ttfamily{*}}\ttfamily{ }\textcolor[HTML]{118dc3}{\ttfamily{x}}\textcolor[HTML]{9a77cf}{\ttfamily{(}}\textcolor[HTML]{e05661}{\ttfamily{i}}\ttfamily{ }\textcolor[HTML]{56b6c2}{\ttfamily{+}}\ttfamily{ }\textcolor[HTML]{ee9025}{\ttfamily{1}}\textcolor[HTML]{56b6c2}{\ttfamily{:}}\textcolor[HTML]{e05661}{\ttfamily{n}}\textcolor[HTML]{9a77cf}{\ttfamily{))}}\textcolor[HTML]{6a6a6a}{\ttfamily{;}}\newline
\hspace*{4ex}\textcolor[HTML]{9a77cf}{\ttfamily{end}}\newline
\hspace*{4ex}\textcolor[HTML]{118dc3}{\ttfamily{x}}\textcolor[HTML]{9a77cf}{\ttfamily{(}}\textcolor[HTML]{e05661}{\ttfamily{n}}\textcolor[HTML]{9a77cf}{\ttfamily{)}}\ttfamily{ }\textcolor[HTML]{56b6c2}{\ttfamily{=}}\ttfamily{ }\textcolor[HTML]{ee9025}{\ttfamily{1.0}}\ttfamily{ }\textcolor[HTML]{56b6c2}{\ttfamily{/}}\ttfamily{ }\textcolor[HTML]{118dc3}{\ttfamily{A}}\textcolor[HTML]{9a77cf}{\ttfamily{(}}\textcolor[HTML]{e05661}{\ttfamily{n}}\textcolor[HTML]{6a6a6a}{\ttfamily{,}}\ttfamily{ }\textcolor[HTML]{e05661}{\ttfamily{n}}\textcolor[HTML]{9a77cf}{\ttfamily{)}}\ttfamily{ }\textcolor[HTML]{56b6c2}{\ttfamily{*}}\ttfamily{ }\textcolor[HTML]{9a77cf}{\ttfamily{(}}\textcolor[HTML]{118dc3}{\ttfamily{b}}\textcolor[HTML]{9a77cf}{\ttfamily{(}}\textcolor[HTML]{e05661}{\ttfamily{n}}\textcolor[HTML]{9a77cf}{\ttfamily{)}}\ttfamily{ }\textcolor[HTML]{56b6c2}{\ttfamily{-}}\ttfamily{ }\textcolor[HTML]{118dc3}{\ttfamily{A}}\textcolor[HTML]{9a77cf}{\ttfamily{(}}\textcolor[HTML]{e05661}{\ttfamily{n}}\textcolor[HTML]{6a6a6a}{\ttfamily{,}}\ttfamily{ }\textcolor[HTML]{ee9025}{\ttfamily{1}}\textcolor[HTML]{56b6c2}{\ttfamily{:}}\textcolor[HTML]{e05661}{\ttfamily{n}}\ttfamily{ }\textcolor[HTML]{56b6c2}{\ttfamily{-}}\ttfamily{ }\textcolor[HTML]{ee9025}{\ttfamily{1}}\textcolor[HTML]{9a77cf}{\ttfamily{)}}\ttfamily{ }\textcolor[HTML]{56b6c2}{\ttfamily{*}}\ttfamily{ }\textcolor[HTML]{118dc3}{\ttfamily{x}}\textcolor[HTML]{9a77cf}{\ttfamily{(}}\textcolor[HTML]{ee9025}{\ttfamily{1}}\textcolor[HTML]{56b6c2}{\ttfamily{:}}\textcolor[HTML]{e05661}{\ttfamily{n}}\ttfamily{ }\textcolor[HTML]{56b6c2}{\ttfamily{-}}\ttfamily{ }\textcolor[HTML]{ee9025}{\ttfamily{1}}\textcolor[HTML]{9a77cf}{\ttfamily{))}}\textcolor[HTML]{6a6a6a}{\ttfamily{;}}\newline
\hspace*{4ex}\textcolor[HTML]{e05661}{\ttfamily{relresid}}\ttfamily{ }\textcolor[HTML]{56b6c2}{\ttfamily{=}}\ttfamily{ }\textcolor[HTML]{118dc3}{\ttfamily{norm}}\textcolor[HTML]{9a77cf}{\ttfamily{(}}\textcolor[HTML]{e05661}{\ttfamily{b}}\ttfamily{ }\textcolor[HTML]{56b6c2}{\ttfamily{-}}\ttfamily{ }\textcolor[HTML]{e05661}{\ttfamily{A}}\ttfamily{ }\textcolor[HTML]{56b6c2}{\ttfamily{*}}\ttfamily{ }\textcolor[HTML]{e05661}{\ttfamily{x}}\textcolor[HTML]{9a77cf}{\ttfamily{)}}\ttfamily{ }\textcolor[HTML]{56b6c2}{\ttfamily{/}}\ttfamily{ }\textcolor[HTML]{118dc3}{\ttfamily{norm}}\textcolor[HTML]{9a77cf}{\ttfamily{(}}\textcolor[HTML]{e05661}{\ttfamily{b}}\textcolor[HTML]{9a77cf}{\ttfamily{)}}\textcolor[HTML]{6a6a6a}{\ttfamily{;}}\newline
\hspace*{4ex}\textcolor[HTML]{9a77cf}{\ttfamily{if}}\ttfamily{ }\textcolor[HTML]{e05661}{\ttfamily{relresid}}\ttfamily{ }\textcolor[HTML]{56b6c2}{\ttfamily{<}}\ttfamily{ }\textcolor[HTML]{e05661}{\ttfamily{tol}}\newline
\hspace*{6ex}\textcolor[HTML]{e05661}{\ttfamily{iter}}\ttfamily{ }\textcolor[HTML]{56b6c2}{\ttfamily{=}}\ttfamily{ }\textcolor[HTML]{e05661}{\ttfamily{k}}\textcolor[HTML]{6a6a6a}{\ttfamily{;}}\newline
\hspace*{6ex}\textcolor[HTML]{9a77cf}{\ttfamily{return}}\newline
\hspace*{4ex}\textcolor[HTML]{9a77cf}{\ttfamily{end}}\newline
\hspace*{2ex}\textcolor[HTML]{9a77cf}{\ttfamily{end}}\newline
\hspace*{0ex}\newline
\hspace*{2ex}\textcolor[HTML]{e05661}{\ttfamily{iter}}\ttfamily{ }\textcolor[HTML]{56b6c2}{\ttfamily{=}}\ttfamily{ }\textcolor[HTML]{56b6c2}{\ttfamily{-}}\textcolor[HTML]{ee9025}{\ttfamily{1}}\textcolor[HTML]{6a6a6a}{\ttfamily{;}}
}
\newpage
\section{Problem20.19: Comparison of the results}
The numerical results are listed below (all logarithms are natural logarithms, i.e., base $e$):
\begin{itemize}
  \item \textbf{Jacobi:} converges in 30 iterations, with log-residual $-33.029204$.
  \item \textbf{Gauss--Seidel:} converges in 21 iterations, with log-residual $-33.886495$.
  \item \textbf{SOR$(\omega=1.1)$:} converges in 21 iterations, with log-residual $-33.365732$.
  \item \textbf{SOR$(\omega=1.2)$:} converges in 28 iterations, with log-residual $-33.997858$.
  \item \textbf{SOR$(\omega=1.3)$:} converges in 35 iterations, with log-residual $-33.057573$.
  \item \textbf{SOR$(\omega=1.5)$:} converges in 60 iterations, with log-residual $-33.292912$.
  \item \textbf{SOR$(\omega=1.9)$:} does not converge.
  \item \textbf{$B\backslash c$:} log-residual $-37.450359$.
\end{itemize}

From these results, we see that the Gauss--Seidel method converges faster than the Jacobi method, since it requires only 21 iterations instead of 30. For the SOR method, the choice of $\omega$ has a strong effect on performance. Among the tested values, $\omega=1.1$ works best in terms of convergence speed, since it converges in 21 iterations, the same as Gauss--Seidel. As $\omega$ increases, the convergence becomes slower, and when $\omega=1.9$, the method does not converge.

In terms of accuracy, the direct solver $B\backslash c$ gives the smallest residual, since its log-residual is $-37.450359$, which is smaller than those of all iterative methods. Therefore, $B\backslash c$ produces the most accurate solution overall.

In conclusion, among the iterative methods, Gauss--Seidel and SOR with $\omega=1.1$ perform best, while the MATLAB command $B\backslash c$ gives the highest accuracy.
\paragraph{Codes}
\subsection{Codes for problem 20.19}
% \lstinputlisting[language=Matlab]{home1_mat/problem20_19.m}
{\noindent\hspace*{0ex}\textcolor[HTML]{e05661}{\ttfamily{n}}\ttfamily{ }\textcolor[HTML]{56b6c2}{\ttfamily{=}}\ttfamily{ }\textcolor[HTML]{ee9025}{\ttfamily{100}}\textcolor[HTML]{6a6a6a}{\ttfamily{;}}\newline
\hspace*{0ex}\textcolor[HTML]{e05661}{\ttfamily{r}}\ttfamily{ }\textcolor[HTML]{56b6c2}{\ttfamily{=}}\ttfamily{ }\textcolor[HTML]{ee9025}{\ttfamily{0.25}}\textcolor[HTML]{6a6a6a}{\ttfamily{;}}\newline
\hspace*{0ex}\textcolor[HTML]{e05661}{\ttfamily{a}}\ttfamily{ }\textcolor[HTML]{56b6c2}{\ttfamily{=}}\ttfamily{ }\textcolor[HTML]{9a77cf}{\ttfamily{(}}\textcolor[HTML]{56b6c2}{\ttfamily{-}}\textcolor[HTML]{e05661}{\ttfamily{r}}\textcolor[HTML]{9a77cf}{\ttfamily{)}}\ttfamily{ }\textcolor[HTML]{56b6c2}{\ttfamily{*}}\ttfamily{ }\textcolor[HTML]{118dc3}{\ttfamily{ones}}\textcolor[HTML]{9a77cf}{\ttfamily{(}}\textcolor[HTML]{e05661}{\ttfamily{n}}\ttfamily{ }\textcolor[HTML]{56b6c2}{\ttfamily{-}}\ttfamily{ }\textcolor[HTML]{ee9025}{\ttfamily{1}}\textcolor[HTML]{6a6a6a}{\ttfamily{,}}\ttfamily{ }\textcolor[HTML]{ee9025}{\ttfamily{1}}\textcolor[HTML]{9a77cf}{\ttfamily{)}}\textcolor[HTML]{6a6a6a}{\ttfamily{;}}\ttfamily{  }\textcolor[HTML]{9b9fa6}{\ttfamily{\% 下对角}}\newline
\hspace*{0ex}\textcolor[HTML]{e05661}{\ttfamily{b}}\ttfamily{ }\textcolor[HTML]{56b6c2}{\ttfamily{=}}\ttfamily{  }\textcolor[HTML]{9a77cf}{\ttfamily{(}}\textcolor[HTML]{ee9025}{\ttfamily{1}}\ttfamily{ }\textcolor[HTML]{56b6c2}{\ttfamily{+}}\ttfamily{ }\textcolor[HTML]{ee9025}{\ttfamily{2}}\ttfamily{ }\textcolor[HTML]{56b6c2}{\ttfamily{*}}\ttfamily{ }\textcolor[HTML]{e05661}{\ttfamily{r}}\textcolor[HTML]{9a77cf}{\ttfamily{)}}\ttfamily{ }\textcolor[HTML]{56b6c2}{\ttfamily{*}}\ttfamily{ }\textcolor[HTML]{118dc3}{\ttfamily{ones}}\textcolor[HTML]{9a77cf}{\ttfamily{(}}\textcolor[HTML]{e05661}{\ttfamily{n}}\textcolor[HTML]{6a6a6a}{\ttfamily{,}}\ttfamily{ }\textcolor[HTML]{ee9025}{\ttfamily{1}}\textcolor[HTML]{9a77cf}{\ttfamily{)}}\textcolor[HTML]{6a6a6a}{\ttfamily{;}}\ttfamily{    }\textcolor[HTML]{9b9fa6}{\ttfamily{\% 主对角}}\newline
\hspace*{0ex}\textcolor[HTML]{e05661}{\ttfamily{c}}\ttfamily{ }\textcolor[HTML]{56b6c2}{\ttfamily{=}}\ttfamily{ }\textcolor[HTML]{9a77cf}{\ttfamily{(}}\textcolor[HTML]{56b6c2}{\ttfamily{-}}\textcolor[HTML]{e05661}{\ttfamily{r}}\textcolor[HTML]{9a77cf}{\ttfamily{)}}\ttfamily{ }\textcolor[HTML]{56b6c2}{\ttfamily{*}}\ttfamily{ }\textcolor[HTML]{118dc3}{\ttfamily{ones}}\textcolor[HTML]{9a77cf}{\ttfamily{(}}\textcolor[HTML]{e05661}{\ttfamily{n}}\ttfamily{ }\textcolor[HTML]{56b6c2}{\ttfamily{-}}\ttfamily{ }\textcolor[HTML]{ee9025}{\ttfamily{1}}\textcolor[HTML]{6a6a6a}{\ttfamily{,}}\ttfamily{ }\textcolor[HTML]{ee9025}{\ttfamily{1}}\textcolor[HTML]{9a77cf}{\ttfamily{)}}\textcolor[HTML]{6a6a6a}{\ttfamily{;}}\ttfamily{  }\textcolor[HTML]{9b9fa6}{\ttfamily{\% 上对角}}\newline
\hspace*{0ex}\newline
\hspace*{0ex}\textcolor[HTML]{e05661}{\ttfamily{A}}\ttfamily{ }\textcolor[HTML]{56b6c2}{\ttfamily{=}}\ttfamily{ }\textcolor[HTML]{118dc3}{\ttfamily{diag}}\textcolor[HTML]{9a77cf}{\ttfamily{(}}\textcolor[HTML]{e05661}{\ttfamily{b}}\textcolor[HTML]{9a77cf}{\ttfamily{)}}\ttfamily{ }\textcolor[HTML]{56b6c2}{\ttfamily{+}}\ttfamily{ }\textcolor[HTML]{118dc3}{\ttfamily{diag}}\textcolor[HTML]{9a77cf}{\ttfamily{(}}\textcolor[HTML]{e05661}{\ttfamily{a}}\textcolor[HTML]{6a6a6a}{\ttfamily{,}}\ttfamily{ }\textcolor[HTML]{56b6c2}{\ttfamily{-}}\textcolor[HTML]{ee9025}{\ttfamily{1}}\textcolor[HTML]{9a77cf}{\ttfamily{)}}\ttfamily{ }\textcolor[HTML]{56b6c2}{\ttfamily{+}}\ttfamily{ }\textcolor[HTML]{118dc3}{\ttfamily{diag}}\textcolor[HTML]{9a77cf}{\ttfamily{(}}\textcolor[HTML]{e05661}{\ttfamily{c}}\textcolor[HTML]{6a6a6a}{\ttfamily{,}}\ttfamily{ }\textcolor[HTML]{ee9025}{\ttfamily{1}}\textcolor[HTML]{9a77cf}{\ttfamily{)}}\textcolor[HTML]{6a6a6a}{\ttfamily{;}}\newline
\hspace*{0ex}\textcolor[HTML]{e05661}{\ttfamily{b}}\ttfamily{ }\textcolor[HTML]{56b6c2}{\ttfamily{=}}\ttfamily{ }\textcolor[HTML]{118dc3}{\ttfamily{ones}}\textcolor[HTML]{9a77cf}{\ttfamily{(}}\textcolor[HTML]{e05661}{\ttfamily{n}}\textcolor[HTML]{6a6a6a}{\ttfamily{,}}\ttfamily{ }\textcolor[HTML]{ee9025}{\ttfamily{1}}\textcolor[HTML]{9a77cf}{\ttfamily{)}}\textcolor[HTML]{6a6a6a}{\ttfamily{;}}\newline
\hspace*{0ex}\textcolor[HTML]{e05661}{\ttfamily{tol}}\ttfamily{ }\textcolor[HTML]{56b6c2}{\ttfamily{=}}\ttfamily{ }\textcolor[HTML]{ee9025}{\ttfamily{0.5e-14}}\textcolor[HTML]{6a6a6a}{\ttfamily{;}}\newline
\hspace*{0ex}\textcolor[HTML]{e05661}{\ttfamily{numiter}}\ttfamily{ }\textcolor[HTML]{56b6c2}{\ttfamily{=}}\ttfamily{ }\textcolor[HTML]{ee9025}{\ttfamily{100}}\textcolor[HTML]{6a6a6a}{\ttfamily{;}}\newline
\hspace*{0ex}\textcolor[HTML]{e05661}{\ttfamily{x0}}\ttfamily{ }\textcolor[HTML]{56b6c2}{\ttfamily{=}}\ttfamily{ }\textcolor[HTML]{118dc3}{\ttfamily{zeros}}\textcolor[HTML]{9a77cf}{\ttfamily{(}}\textcolor[HTML]{e05661}{\ttfamily{n}}\textcolor[HTML]{6a6a6a}{\ttfamily{,}}\ttfamily{ }\textcolor[HTML]{ee9025}{\ttfamily{1}}\textcolor[HTML]{9a77cf}{\ttfamily{)}}\textcolor[HTML]{6a6a6a}{\ttfamily{;}}\newline
\hspace*{0ex}\newline
\hspace*{0ex}\textcolor[HTML]{9a77cf}{\ttfamily{[}}\textcolor[HTML]{9a77cf}{\ttfamily{\textasciitilde{}}}\textcolor[HTML]{6a6a6a}{\ttfamily{,}}\ttfamily{ }\textcolor[HTML]{eea825}{\ttfamily{GSiter}}\textcolor[HTML]{6a6a6a}{\ttfamily{,}}\ttfamily{ }\textcolor[HTML]{eea825}{\ttfamily{GSresider}}\textcolor[HTML]{9a77cf}{\ttfamily{]}}\ttfamily{ }\textcolor[HTML]{56b6c2}{\ttfamily{=}}\ttfamily{ }\textcolor[HTML]{118dc3}{\ttfamily{gausseidel}}\textcolor[HTML]{9a77cf}{\ttfamily{(}}\textcolor[HTML]{e05661}{\ttfamily{A}}\textcolor[HTML]{6a6a6a}{\ttfamily{,}}\ttfamily{ }\textcolor[HTML]{e05661}{\ttfamily{b}}\textcolor[HTML]{6a6a6a}{\ttfamily{,}}\ttfamily{ }\textcolor[HTML]{e05661}{\ttfamily{x0}}\textcolor[HTML]{6a6a6a}{\ttfamily{,}}\ttfamily{ }\textcolor[HTML]{e05661}{\ttfamily{tol}}\textcolor[HTML]{6a6a6a}{\ttfamily{,}}\ttfamily{ }\textcolor[HTML]{e05661}{\ttfamily{numiter}}\textcolor[HTML]{9a77cf}{\ttfamily{)}}\textcolor[HTML]{6a6a6a}{\ttfamily{;}}\newline
\hspace*{0ex}\textcolor[HTML]{eea825}{\ttfamily{GSlogre}}\ttfamily{ }\textcolor[HTML]{56b6c2}{\ttfamily{=}}\ttfamily{ }\textcolor[HTML]{118dc3}{\ttfamily{log}}\textcolor[HTML]{9a77cf}{\ttfamily{(}}\textcolor[HTML]{eea825}{\ttfamily{GSresider}}\textcolor[HTML]{9a77cf}{\ttfamily{)}}\textcolor[HTML]{6a6a6a}{\ttfamily{;}}\newline
\hspace*{0ex}\textcolor[HTML]{9a77cf}{\ttfamily{if}}\ttfamily{ }\textcolor[HTML]{eea825}{\ttfamily{GSiter}}\ttfamily{ }\textcolor[HTML]{56b6c2}{\ttfamily{>}}\ttfamily{ }\textcolor[HTML]{ee9025}{\ttfamily{0}}\newline
\hspace*{2ex}\textcolor[HTML]{118dc3}{\ttfamily{fprintf}}\textcolor[HTML]{9a77cf}{\ttfamily{(}}\textcolor[HTML]{56b6c2}{\ttfamily{'}}\textcolor[HTML]{1da912}{\ttfamily{Gauss-Seidel converges at step }}\textcolor[HTML]{118dc3}{\ttfamily{\%f}}\textcolor[HTML]{1da912}{\ttfamily{ with residual }}\textcolor[HTML]{118dc3}{\ttfamily{\%f}}\textcolor[HTML]{1da912}{\ttfamily{ and log-residual }}\textcolor[HTML]{118dc3}{\ttfamily{\%f}}\textcolor[HTML]{56b6c2}{\ttfamily{\textbackslash{}n}}\textcolor[HTML]{56b6c2}{\ttfamily{'}}\textcolor[HTML]{6a6a6a}{\ttfamily{,}}\ttfamily{ }\textcolor[HTML]{9b9fa6}{\ttfamily{...}}\newline
\hspace*{10ex}\textcolor[HTML]{eea825}{\ttfamily{GSiter}}\textcolor[HTML]{6a6a6a}{\ttfamily{,}}\ttfamily{ }\textcolor[HTML]{eea825}{\ttfamily{GSresider}}\textcolor[HTML]{6a6a6a}{\ttfamily{,}}\ttfamily{ }\textcolor[HTML]{eea825}{\ttfamily{GSlogre}}\textcolor[HTML]{9a77cf}{\ttfamily{)}}\textcolor[HTML]{6a6a6a}{\ttfamily{;}}\newline
\hspace*{0ex}\textcolor[HTML]{9a77cf}{\ttfamily{else}}\newline
\hspace*{2ex}\textcolor[HTML]{118dc3}{\ttfamily{fprintf}}\textcolor[HTML]{9a77cf}{\ttfamily{(}}\textcolor[HTML]{56b6c2}{\ttfamily{'}}\textcolor[HTML]{1da912}{\ttfamily{Gauss-Seidel does not converge}}\textcolor[HTML]{56b6c2}{\ttfamily{\textbackslash{}n}}\textcolor[HTML]{56b6c2}{\ttfamily{'}}\textcolor[HTML]{9a77cf}{\ttfamily{)}}\textcolor[HTML]{6a6a6a}{\ttfamily{;}}\newline
\hspace*{0ex}\textcolor[HTML]{9a77cf}{\ttfamily{end}}\newline
\hspace*{0ex}\textcolor[HTML]{9a77cf}{\ttfamily{[}}\textcolor[HTML]{9a77cf}{\ttfamily{\textasciitilde{}}}\textcolor[HTML]{6a6a6a}{\ttfamily{,}}\ttfamily{ }\textcolor[HTML]{eea825}{\ttfamily{Jiter}}\textcolor[HTML]{6a6a6a}{\ttfamily{,}}\ttfamily{ }\textcolor[HTML]{eea825}{\ttfamily{Jresider}}\textcolor[HTML]{9a77cf}{\ttfamily{]}}\ttfamily{ }\textcolor[HTML]{56b6c2}{\ttfamily{=}}\ttfamily{ }\textcolor[HTML]{118dc3}{\ttfamily{Jacobi}}\textcolor[HTML]{9a77cf}{\ttfamily{(}}\textcolor[HTML]{e05661}{\ttfamily{A}}\textcolor[HTML]{6a6a6a}{\ttfamily{,}}\ttfamily{ }\textcolor[HTML]{e05661}{\ttfamily{b}}\textcolor[HTML]{6a6a6a}{\ttfamily{,}}\ttfamily{ }\textcolor[HTML]{e05661}{\ttfamily{x0}}\textcolor[HTML]{6a6a6a}{\ttfamily{,}}\ttfamily{ }\textcolor[HTML]{e05661}{\ttfamily{tol}}\textcolor[HTML]{6a6a6a}{\ttfamily{,}}\ttfamily{ }\textcolor[HTML]{e05661}{\ttfamily{numiter}}\textcolor[HTML]{9a77cf}{\ttfamily{)}}\textcolor[HTML]{6a6a6a}{\ttfamily{;}}\newline
\hspace*{0ex}\textcolor[HTML]{eea825}{\ttfamily{Jlogre}}\ttfamily{ }\textcolor[HTML]{56b6c2}{\ttfamily{=}}\ttfamily{ }\textcolor[HTML]{118dc3}{\ttfamily{log}}\textcolor[HTML]{9a77cf}{\ttfamily{(}}\textcolor[HTML]{eea825}{\ttfamily{Jresider}}\textcolor[HTML]{9a77cf}{\ttfamily{)}}\textcolor[HTML]{6a6a6a}{\ttfamily{;}}\newline
\hspace*{0ex}\textcolor[HTML]{9a77cf}{\ttfamily{if}}\ttfamily{ }\textcolor[HTML]{eea825}{\ttfamily{Jiter}}\ttfamily{ }\textcolor[HTML]{56b6c2}{\ttfamily{>}}\ttfamily{ }\textcolor[HTML]{ee9025}{\ttfamily{0}}\newline
\hspace*{2ex}\textcolor[HTML]{118dc3}{\ttfamily{fprintf}}\textcolor[HTML]{9a77cf}{\ttfamily{(}}\textcolor[HTML]{56b6c2}{\ttfamily{'}}\textcolor[HTML]{1da912}{\ttfamily{Jacobi converges at step }}\textcolor[HTML]{118dc3}{\ttfamily{\%f}}\textcolor[HTML]{1da912}{\ttfamily{ with residual }}\textcolor[HTML]{118dc3}{\ttfamily{\%f}}\textcolor[HTML]{1da912}{\ttfamily{ and log-residual }}\textcolor[HTML]{118dc3}{\ttfamily{\%f}}\textcolor[HTML]{56b6c2}{\ttfamily{\textbackslash{}n}}\textcolor[HTML]{56b6c2}{\ttfamily{'}}\textcolor[HTML]{6a6a6a}{\ttfamily{,}}\ttfamily{ }\textcolor[HTML]{9b9fa6}{\ttfamily{...}}\newline
\hspace*{10ex}\textcolor[HTML]{eea825}{\ttfamily{Jiter}}\textcolor[HTML]{6a6a6a}{\ttfamily{,}}\ttfamily{ }\textcolor[HTML]{eea825}{\ttfamily{Jresider}}\textcolor[HTML]{6a6a6a}{\ttfamily{,}}\ttfamily{ }\textcolor[HTML]{eea825}{\ttfamily{Jlogre}}\textcolor[HTML]{9a77cf}{\ttfamily{)}}\textcolor[HTML]{6a6a6a}{\ttfamily{;}}\newline
\hspace*{0ex}\newline
\hspace*{0ex}\textcolor[HTML]{9a77cf}{\ttfamily{else}}\newline
\hspace*{2ex}\textcolor[HTML]{118dc3}{\ttfamily{fprintf}}\textcolor[HTML]{9a77cf}{\ttfamily{(}}\textcolor[HTML]{56b6c2}{\ttfamily{'}}\textcolor[HTML]{1da912}{\ttfamily{Jacobi does not converge}}\textcolor[HTML]{56b6c2}{\ttfamily{\textbackslash{}n}}\textcolor[HTML]{56b6c2}{\ttfamily{'}}\textcolor[HTML]{9a77cf}{\ttfamily{)}}\textcolor[HTML]{6a6a6a}{\ttfamily{;}}\newline
\hspace*{0ex}\textcolor[HTML]{9a77cf}{\ttfamily{end}}\newline
\hspace*{0ex}\textcolor[HTML]{e05661}{\ttfamily{W}}\ttfamily{ }\textcolor[HTML]{56b6c2}{\ttfamily{=}}\ttfamily{ }\textcolor[HTML]{9a77cf}{\ttfamily{[}}\textcolor[HTML]{ee9025}{\ttfamily{1.1}}\textcolor[HTML]{6a6a6a}{\ttfamily{,}}\ttfamily{ }\textcolor[HTML]{ee9025}{\ttfamily{1.2}}\textcolor[HTML]{6a6a6a}{\ttfamily{,}}\ttfamily{ }\textcolor[HTML]{ee9025}{\ttfamily{1.3}}\textcolor[HTML]{6a6a6a}{\ttfamily{,}}\ttfamily{ }\textcolor[HTML]{ee9025}{\ttfamily{1.5}}\textcolor[HTML]{6a6a6a}{\ttfamily{,}}\ttfamily{ }\textcolor[HTML]{ee9025}{\ttfamily{1.9}}\textcolor[HTML]{9a77cf}{\ttfamily{]}}\textcolor[HTML]{6a6a6a}{\ttfamily{;}}\newline
\hspace*{0ex}\textcolor[HTML]{9a77cf}{\ttfamily{for}}\ttfamily{ }\textcolor[HTML]{e05661}{\ttfamily{w}}\ttfamily{ }\textcolor[HTML]{56b6c2}{\ttfamily{=}}\ttfamily{ }\textcolor[HTML]{e05661}{\ttfamily{W}}\newline
\hspace*{2ex}\textcolor[HTML]{9a77cf}{\ttfamily{[}}\textcolor[HTML]{9a77cf}{\ttfamily{\textasciitilde{}}}\textcolor[HTML]{6a6a6a}{\ttfamily{,}}\ttfamily{ }\textcolor[HTML]{eea825}{\ttfamily{Soriter}}\textcolor[HTML]{6a6a6a}{\ttfamily{,}}\ttfamily{ }\textcolor[HTML]{eea825}{\ttfamily{Sorresider}}\textcolor[HTML]{9a77cf}{\ttfamily{]}}\ttfamily{ }\textcolor[HTML]{56b6c2}{\ttfamily{=}}\ttfamily{ }\textcolor[HTML]{118dc3}{\ttfamily{sor}}\textcolor[HTML]{9a77cf}{\ttfamily{(}}\textcolor[HTML]{e05661}{\ttfamily{A}}\textcolor[HTML]{6a6a6a}{\ttfamily{,}}\ttfamily{ }\textcolor[HTML]{e05661}{\ttfamily{b}}\textcolor[HTML]{6a6a6a}{\ttfamily{,}}\ttfamily{ }\textcolor[HTML]{e05661}{\ttfamily{x0}}\textcolor[HTML]{6a6a6a}{\ttfamily{,}}\ttfamily{ }\textcolor[HTML]{e05661}{\ttfamily{w}}\textcolor[HTML]{6a6a6a}{\ttfamily{,}}\ttfamily{ }\textcolor[HTML]{e05661}{\ttfamily{tol}}\textcolor[HTML]{6a6a6a}{\ttfamily{,}}\ttfamily{ }\textcolor[HTML]{e05661}{\ttfamily{numiter}}\textcolor[HTML]{9a77cf}{\ttfamily{)}}\textcolor[HTML]{6a6a6a}{\ttfamily{;}}\newline
\hspace*{2ex}\textcolor[HTML]{eea825}{\ttfamily{Sorlogre}}\ttfamily{ }\textcolor[HTML]{56b6c2}{\ttfamily{=}}\ttfamily{ }\textcolor[HTML]{118dc3}{\ttfamily{log}}\textcolor[HTML]{9a77cf}{\ttfamily{(}}\textcolor[HTML]{eea825}{\ttfamily{Sorresider}}\textcolor[HTML]{9a77cf}{\ttfamily{)}}\textcolor[HTML]{6a6a6a}{\ttfamily{;}}\newline
\hspace*{2ex}\textcolor[HTML]{9a77cf}{\ttfamily{if}}\ttfamily{ }\textcolor[HTML]{eea825}{\ttfamily{Soriter}}\ttfamily{ }\textcolor[HTML]{56b6c2}{\ttfamily{>}}\ttfamily{ }\textcolor[HTML]{ee9025}{\ttfamily{0}}\newline
\hspace*{4ex}\textcolor[HTML]{118dc3}{\ttfamily{fprintf}}\textcolor[HTML]{9a77cf}{\ttfamily{(}}\textcolor[HTML]{56b6c2}{\ttfamily{'}}\textcolor[HTML]{1da912}{\ttfamily{SOR(}}\textcolor[HTML]{118dc3}{\ttfamily{\%.1f}}\textcolor[HTML]{1da912}{\ttfamily{) converges at step }}\textcolor[HTML]{118dc3}{\ttfamily{\%f}}\textcolor[HTML]{1da912}{\ttfamily{ with residual }}\textcolor[HTML]{118dc3}{\ttfamily{\%f}}\textcolor[HTML]{1da912}{\ttfamily{ and log-residual }}\textcolor[HTML]{118dc3}{\ttfamily{\%f}}\textcolor[HTML]{56b6c2}{\ttfamily{\textbackslash{}n}}\textcolor[HTML]{56b6c2}{\ttfamily{'}}\textcolor[HTML]{6a6a6a}{\ttfamily{,}}\ttfamily{ }\textcolor[HTML]{9b9fa6}{\ttfamily{...}}\newline
\hspace*{12ex}\textcolor[HTML]{e05661}{\ttfamily{w}}\textcolor[HTML]{6a6a6a}{\ttfamily{,}}\ttfamily{ }\textcolor[HTML]{eea825}{\ttfamily{Soriter}}\textcolor[HTML]{6a6a6a}{\ttfamily{,}}\ttfamily{ }\textcolor[HTML]{eea825}{\ttfamily{Sorresider}}\textcolor[HTML]{6a6a6a}{\ttfamily{,}}\ttfamily{ }\textcolor[HTML]{eea825}{\ttfamily{Sorlogre}}\textcolor[HTML]{9a77cf}{\ttfamily{)}}\textcolor[HTML]{6a6a6a}{\ttfamily{;}}\newline
\hspace*{2ex}\textcolor[HTML]{9a77cf}{\ttfamily{else}}\newline
\hspace*{4ex}\textcolor[HTML]{118dc3}{\ttfamily{fprintf}}\textcolor[HTML]{9a77cf}{\ttfamily{(}}\textcolor[HTML]{56b6c2}{\ttfamily{'}}\textcolor[HTML]{1da912}{\ttfamily{SOR(}}\textcolor[HTML]{118dc3}{\ttfamily{\%.1f}}\textcolor[HTML]{1da912}{\ttfamily{) does not converge}}\textcolor[HTML]{56b6c2}{\ttfamily{\textbackslash{}n}}\textcolor[HTML]{56b6c2}{\ttfamily{'}}\textcolor[HTML]{6a6a6a}{\ttfamily{,}}\ttfamily{ }\textcolor[HTML]{e05661}{\ttfamily{w}}\textcolor[HTML]{9a77cf}{\ttfamily{)}}\textcolor[HTML]{6a6a6a}{\ttfamily{;}}\newline
\hspace*{2ex}\textcolor[HTML]{9a77cf}{\ttfamily{end}}\newline
\hspace*{0ex}\textcolor[HTML]{9a77cf}{\ttfamily{end}}\newline
\hspace*{0ex}\textcolor[HTML]{e05661}{\ttfamily{x\_dir}}\ttfamily{ }\textcolor[HTML]{56b6c2}{\ttfamily{=}}\ttfamily{ }\textcolor[HTML]{e05661}{\ttfamily{A}}\ttfamily{ }\textcolor[HTML]{56b6c2}{\ttfamily{\textbackslash{}}}\ttfamily{ }\textcolor[HTML]{e05661}{\ttfamily{b}}\textcolor[HTML]{6a6a6a}{\ttfamily{;}}\newline
\hspace*{0ex}\textcolor[HTML]{e05661}{\ttfamily{x\_dir\_relresid}}\ttfamily{ }\textcolor[HTML]{56b6c2}{\ttfamily{=}}\ttfamily{ }\textcolor[HTML]{118dc3}{\ttfamily{norm}}\textcolor[HTML]{9a77cf}{\ttfamily{(}}\textcolor[HTML]{e05661}{\ttfamily{b}}\ttfamily{ }\textcolor[HTML]{56b6c2}{\ttfamily{-}}\ttfamily{ }\textcolor[HTML]{e05661}{\ttfamily{A}}\ttfamily{ }\textcolor[HTML]{56b6c2}{\ttfamily{*}}\ttfamily{ }\textcolor[HTML]{e05661}{\ttfamily{x\_dir}}\textcolor[HTML]{9a77cf}{\ttfamily{)}}\ttfamily{ }\textcolor[HTML]{56b6c2}{\ttfamily{/}}\ttfamily{ }\textcolor[HTML]{118dc3}{\ttfamily{norm}}\textcolor[HTML]{9a77cf}{\ttfamily{(}}\textcolor[HTML]{e05661}{\ttfamily{b}}\textcolor[HTML]{9a77cf}{\ttfamily{)}}\textcolor[HTML]{6a6a6a}{\ttfamily{;}}\newline
\hspace*{0ex}\textcolor[HTML]{e05661}{\ttfamily{log\_x\_dir\_residual}}\ttfamily{ }\textcolor[HTML]{56b6c2}{\ttfamily{=}}\ttfamily{ }\textcolor[HTML]{118dc3}{\ttfamily{log}}\textcolor[HTML]{9a77cf}{\ttfamily{(}}\textcolor[HTML]{e05661}{\ttfamily{x\_dir\_relresid}}\textcolor[HTML]{9a77cf}{\ttfamily{)}}\textcolor[HTML]{6a6a6a}{\ttfamily{;}}\newline
\hspace*{0ex}\textcolor[HTML]{118dc3}{\ttfamily{fprintf}}\textcolor[HTML]{9a77cf}{\ttfamily{(}}\textcolor[HTML]{56b6c2}{\ttfamily{'}}\textcolor[HTML]{1da912}{\ttfamily{B}}\textcolor[HTML]{56b6c2}{\ttfamily{\textbackslash{}\textbackslash{}}}\textcolor[HTML]{1da912}{\ttfamily{c residual is }}\textcolor[HTML]{118dc3}{\ttfamily{\%f}}\textcolor[HTML]{1da912}{\ttfamily{ and log-residual }}\textcolor[HTML]{118dc3}{\ttfamily{\%f}}\textcolor[HTML]{56b6c2}{\ttfamily{\textbackslash{}n}}\textcolor[HTML]{56b6c2}{\ttfamily{'}}\textcolor[HTML]{6a6a6a}{\ttfamily{,}}\ttfamily{ }\textcolor[HTML]{9b9fa6}{\ttfamily{...}}\newline
\hspace*{8ex}\textcolor[HTML]{e05661}{\ttfamily{x\_dir\_relresid}}\textcolor[HTML]{6a6a6a}{\ttfamily{,}}\ttfamily{ }\textcolor[HTML]{e05661}{\ttfamily{log\_x\_dir\_residual}}\textcolor[HTML]{9a77cf}{\ttfamily{)}}\textcolor[HTML]{6a6a6a}{\ttfamily{;}}
}
\end{document}
